%% 08_conclusion.tex

\section{Conclusion}
\label{sec:conclusion}

We presented the first paired-data comparison of narrative RAG, naive structured FHIR
RAG, and resource-aware structured FHIR RAG on specialty-medication adherence-indicator
question answering, where information content is held constant across representations
by constructing all three retrieval substrates from the same FHIR bundles.
Narrative RAG achieves the highest exact-match accuracy (40.6\% vs.\ 33.4--35.3\%).
A 1024-token sensitivity sweep on a stratified 495-question sample shows the
structured systems gain $+3.8$ to $+5.1$~pp when the answer budget is doubled while
the narrative system is unchanged, narrowing but not eliminating A's lead; the
A $>$ B $>$ C ordering is therefore not a token-budget artefact alone.
Temporal-anchor failures are a structural weakness of the narrative format that
persists independent of budget.  In the feature-extraction arm, structured FHIR
features substantially outperform both narrative-derived and retrieval-trace features
for adherence classification, demonstrating the representation dissociation that is
the paper's central empirical claim: winning natural-language QA and winning downstream
adherence-outcome ML are different objectives, and the same representation does not
simultaneously optimise both.

Code (three retrieval systems, evaluation harness, feature extractors) and analysis
scripts are released as a reusable benchmark for the specialty-pharmacy QA research
community; dataset release (200 patients, 13,800 questions, five PSP adherence-indicator
families) under CC-BY-4.0 is pending employer IP clearance.

Future directions include: (1) extending the budget sweep to additional points
(256, 768, 2048) and to the full 13,800-question matrix; (2) evaluating on real
de-identified PSP data; (3) testing robustness across multiple answer LLMs and
embedding models; (4) extending the question bank to additional specialty-medication
classes and real-world edge cases (infusion regimens, variable-interval dosing,
bridge therapies); (5) constructing richer synthetic adherence labels to stress-test
the feature-extraction arm beyond the perfectly-scheduled synthetic scenario; and
(6) a drop-one ablation over System~C's three resource-aware primitives (typed filter,
reference-chain traversal, temporal pre-filter) to disentangle which primitive drives
the recall@$k$ gain over System~B.
